\documentclass{article}

\usepackage{amsmath, amssymb}
\usepackage{graphicx}
\usepackage{geometry}
\usepackage{fancyhdr}
\usepackage{enumitem}
\usepackage[french]{babel}
\usepackage[utf8]{inputenc}
\usepackage[T1]{fontenc}

\pagestyle{fancy}
\fancyhf{}
\lhead{Correction du Contrôle 2 (B)}
\rhead{Première Spécialité Mathématiques}
\cfoot{\thepage}

\begin{document}

\section*{Correction du Contrôle 2 (B)}

\subsection*{Exercice 1 (6 points)}

\textbf{1. Pour chacune des suites suivantes, indiquer si elle est arithmétique ou géométrique, et préciser pourquoi. Justifier si elle ne l'est pas.}

\begin{enumerate}[label=\alph*)]

\item \( u_0 = 2, \quad u_{n+1} = \dfrac{4 u_n}{3} \)

\textbf{Correction :}

Calculons le quotient de deux termes consécutifs :

\[ \dfrac{u_{n+1}}{u_n} = \dfrac{\dfrac{4 u_n}{3}}{u_n} = \dfrac{4}{3} \]

Le quotient est constant, donc la suite est \textbf{géométrique} de raison \( q = \dfrac{4}{3} \).

\bigskip

\item \( u_n = 1{,}5 \times \left( \dfrac{2}{3} \right)^n \)

\textbf{Correction :}

La suite est donnée sous forme explicite avec une puissance de \( q = \dfrac{2}{3} \), donc c'est une suite \textbf{géométrique} de raison \( q = \dfrac{2}{3} \).

\bigskip

\item \( u_1 = 2, \quad u_{n+1} = u_n + \dfrac{5}{3} \)

\textbf{Correction :}

Calculons la différence entre deux termes consécutifs :

\[ u_{n+1} - u_n = \dfrac{5}{3} \]

La différence est constante, donc la suite est \textbf{arithmétique} de raison \( r = \dfrac{5}{3} \).

\bigskip

\item \( u_n = n^2 - 8n \)

\textbf{Correction :}

La différence \( u_{n+1} - u_n \) et le quotient \( \dfrac{u_{n+1}}{u_n} \) dépendent de \( n \) et ne sont pas constants.

Donc, la suite n'est \textbf{ni arithmétique ni géométrique}.

\bigskip

\item \( u_n = 7 - \pi n \)

\textbf{Correction :}

Calculons la différence entre deux termes consécutifs :

\[ u_{n+1} - u_n = [7 - \pi(n+1)] - [7 - \pi n] = -\pi \]

La différence est constante, donc la suite est \textbf{arithmétique} de raison \( r = -\pi \).

\bigskip

\item \( u_0 = 1, \quad u_{n+1} = 3n u_n \)

\textbf{Correction :}

La différence \( u_{n+1} - u_n \) et le quotient \( \dfrac{u_{n+1}}{u_n} \) dépendent de \( n \).

Donc, la suite n'est \textbf{ni arithmétique ni géométrique}.

\end{enumerate}

\bigskip

\textbf{2. Pour chaque suite, déterminer ses variations (croissante, décroissante) sur \( \mathbb{N} \).}

\begin{enumerate}[label=\alph*)]

\item Suite géométrique de raison \( q = \dfrac{4}{3} > 1 \).

Donc, la suite est \textbf{strictement croissante}.

\bigskip

\item Suite géométrique de raison \( q = \dfrac{2}{3} \) avec \( 0 < q < 1 \).

Donc, la suite est \textbf{strictement décroissante}.

\bigskip

\item Suite arithmétique de raison \( r = \dfrac{5}{3} > 0 \).

Donc, la suite est \textbf{strictement croissante}.

\bigskip

\item \( u_n = n^2 - 8n \)

Calculons la différence \( u_{n+1} - u_n \) :

\[ u_{n+1} - u_n = (n+1)^2 - 8(n+1) - [n^2 - 8n] = 2n - 7 \]

- Pour \( n \leq 3 \), \( u_{n+1} - u_n \leq -1 \), la suite est \textbf{décroissante}.
- Pour \( n = 4 \), \( u_{5} - u_{4} = 2 \times 4 - 7 = 1 > 0 \), la suite commence à \textbf{croître}.
- Pour \( n \geq 4 \), \( u_{n+1} - u_n \geq 1 \), la suite est \textbf{croissante}.

Donc, la suite est \textbf{décroissante} pour \( n \leq 3 \), atteint un minimum à \( n = 4 \), puis est \textbf{croissante} pour \( n \geq 4 \).

\bigskip

\item Suite arithmétique de raison \( r = -\pi < 0 \).

Donc, la suite est \textbf{strictement décroissante}.

\bigskip

\item Calculons les premiers termes :

\begin{itemize}
    \item \( u_0 = 1 \)
    \item \( u_1 = 3 \times 0 \times u_0 = 0 \)
    \item Pour \( n \geq 1 \), \( u_n = 0 \) car \( u_{n+1} = 3n u_n = 3n \times 0 = 0 \).
\end{itemize}

La suite décroît de \( u_0 = 1 \) à \( u_1 = 0 \), puis reste \textbf{constante} égale à 0.

\end{enumerate}

\bigskip

\textbf{3. Lorsque c'est possible, exprimer chaque suite de manière explicite et calculer son 100\ieme{} terme.}

\begin{enumerate}[label=\alph*)]

\item Suite géométrique de premier terme \( u_0 = 2 \) et raison \( q = \dfrac{4}{3} \).

La formule explicite est :

\[ u_n = u_0 \times q^n = 2 \times \left( \dfrac{4}{3} \right)^n \]

Calcul du 100\ieme{} terme :

\[ u_{100} = 2 \times \left( \dfrac{4}{3} \right)^{100} \]

\bigskip

\item Suite géométrique donnée explicitement :

\[ u_n = 1{,}5 \times \left( \dfrac{2}{3} \right)^n \]

Calcul du 100\ieme{} terme :

\[ u_{100} = 1{,}5 \times \left( \dfrac{2}{3} \right)^{100} \]

\bigskip

\item Suite arithmétique de premier terme \( u_1 = 2 \) et raison \( r = \dfrac{5}{3} \).

La formule explicite est :

\[ u_n = u_1 + (n - 1) r = 2 + (n - 1) \times \dfrac{5}{3} \]

Calcul du 100\ieme{} terme :

\[ u_{100} = 2 + 99 \times \dfrac{5}{3} = 2 + \dfrac{495}{3} = 2 + 165 = 167 \]

\bigskip

\item La suite est déjà définie explicitement :

\[ u_n = n^2 - 8n \]

Calcul du 100\ieme{} terme :

\[ u_{100} = 100^2 - 8 \times 100 = 10\,000 - 800 = 9\,200 \]

\bigskip

\item La suite est déjà définie explicitement :

\[ u_n = 7 - \pi n \]

Calcul du 100\ieme{} terme :

\[ u_{100} = 7 - 100 \pi \]

\bigskip

\item Il n'est pas possible d'exprimer cette suite de manière explicite simple.

Cependant, comme \( u_n = 0 \) pour \( n \geq 1 \), on a :

\[ u_{100} = 0 \]

\end{enumerate}

\bigskip

\subsection*{Exercice 2 (4 points)}

\textbf{1. Calculer les sommes suivantes :}

\begin{enumerate}[label=\alph*)]

\item \( S = 1 + \dfrac{1}{5} + \dfrac{1}{25} + \dots + \dfrac{1}{15\,625} \)

\textbf{Correction :}

Il s'agit d'une suite géométrique de premier terme \( u_0 = 1 \) et raison \( q = \dfrac{1}{5} \).

Le terme général est :

\[ u_n = u_0 \times q^n = \left( \dfrac{1}{5} \right)^n \]

Trouvons le nombre de termes \( n \) :

Sachant que \( 5^6 = 15\,625 \), donc \( \left( \dfrac{1}{5} \right)^6 = \dfrac{1}{15\,625} \).

Ainsi, \( n = 6 \).

La somme des termes est donnée par :

\[ S = u_0 \times \dfrac{1 - q^{n+1}}{1 - q} \]

\[ S = 1 \times \dfrac{1 - \left( \dfrac{1}{5} \right)^{7}}{1 - \dfrac{1}{5}} = \dfrac{1 - \dfrac{1}{78\,125}}{\dfrac{4}{5}} = \dfrac{5}{4} \left( 1 - \dfrac{1}{78\,125} \right) \]

Comme \( \dfrac{1}{78\,125} \) est très petit, on peut approximer :

\[ S \approx \dfrac{5}{4} \]

\bigskip

\item \( S = 5 + 9 + 13 + \dots + 493 \)

\textbf{Correction :}

Il s'agit d'une suite arithmétique de premier terme \( u_1 = 5 \) et raison \( r = 4 \).

Trouvons le nombre de termes \( n \) :

\[ u_n = u_1 + (n - 1) r \]

\[ 493 = 5 + (n - 1) \times 4 \]

\[ 493 - 5 = 4(n - 1) \]

\[ 488 = 4(n - 1) \]

\[ n - 1 = \dfrac{488}{4} = 122 \]

\[ n = 123 \]

La somme des termes est donnée par :

\[ S = \dfrac{n(u_1 + u_n)}{2} \]

\[ S = \dfrac{123(5 + 493)}{2} = \dfrac{123 \times 498}{2} = 61{,}5 \times 498 = 30\,657 \]

\bigskip

\item \( S = \dfrac{4}{3} + \dfrac{4}{9} + \dfrac{4}{27} + \dots + \dfrac{4}{6\,561} \)

\textbf{Correction :}

Il s'agit d'une suite géométrique de premier terme \( u_1 = \dfrac{4}{3} \) et raison \( q = \dfrac{1}{3} \).

Le terme général est :

\[ u_n = u_1 \times q^{n - 1} \]

Trouvons le nombre de termes \( n \) tel que :

\[ u_n = \dfrac{4}{6\,561} = \dfrac{4}{3} \times \dfrac{1}{2\,187} \]

Sachant que \( 3^7 = 2\,187 \), donc \( \left( \dfrac{1}{3} \right)^{7} = \dfrac{1}{2\,187} \). Ainsi, \( n - 1 = 7 \), donc \( n = 8 \).

La somme des termes est donnée par :

\[ S = u_1 \times \dfrac{1 - q^n}{1 - q} \]

\[ S = \dfrac{4}{3} \times \dfrac{1 - \left( \dfrac{1}{3} \right)^8}{1 - \dfrac{1}{3}} = \dfrac{4}{3} \times \dfrac{1 - \dfrac{1}{6\,561}}{\dfrac{2}{3}} \]

\[ S = 2 \left( 1 - \dfrac{1}{6\,561} \right) \]

Comme \( \dfrac{1}{6\,561} \) est très petit, on peut approximer :

\[ S \approx 2 \]

\end{enumerate}

\bigskip

\textbf{2. Exprimer la somme \( \displaystyle \sum_{k=2}^{10} (6k + 2) \) sous la forme \( a + b + \dots + c \).}

\textbf{Correction :}

Calculons les termes pour \( k \) de 2 à 10 :

\begin{align*}
k = 2 &\Rightarrow 6 \times 2 + 2 = 14 \\
k = 3 &\Rightarrow 6 \times 3 + 2 = 20 \\
k = 4 &\Rightarrow 6 \times 4 + 2 = 26 \\
k = 5 &\Rightarrow 6 \times 5 + 2 = 32 \\
k = 6 &\Rightarrow 6 \times 6 + 2 = 38 \\
k = 7 &\Rightarrow 6 \times 7 + 2 = 44 \\
k = 8 &\Rightarrow 6 \times 8 + 2 = 50 \\
k = 9 &\Rightarrow 6 \times 9 + 2 = 56 \\
k = 10 &\Rightarrow 6 \times 10 + 2 = 62 \\
\end{align*}

Donc, la somme est :

\[ 14 + 20 + 26 + 32 + 38 + 44 + 50 + 56 + 62 \]

\bigskip

\subsection*{Exercice 3 (4 points)}

\textbf{1. Dans chaque cas, on donne deux termes d'une suite arithmétique ou géométrique \( (u_n) \). Déterminer la raison et le premier terme \( u_0 \) de cette suite.}

\begin{enumerate}[label=\alph*)]

\item Suite arithmétique, \( u_7 = 31 \) et \( u_{20} = 45 \)

\textbf{Correction :}

La formule générale d'une suite arithmétique est :

\[ u_n = u_0 + n r \]

On a :

\begin{align*}
u_7 &= u_0 + 7 r = 31 \\
u_{20} &= u_0 + 20 r = 45
\end{align*}

En soustrayant les deux équations :

\begin{alignat}{2}
    && (u_{20} - u_7) &= 45 - 31\\
    \iff && (20 r - 7 r) &= 14\\
    \iff && 13r &= 14
\end{alignat}

Donc :

\[ r = \dfrac{14}{13} \]

Ensuite, on calcule \( u_0 \) :

\[ u_0 = u_7 - 7 r = 31 - 7 \times \dfrac{14}{13} = 31 - \dfrac{98}{13} = \dfrac{403}{13} - \dfrac{98}{13} = \dfrac{305}{13} \]

\bigskip

\item Suite géométrique, \( u_3 = 27 \) et \( u_{6} = 729 \)

\textbf{Correction :}

La formule générale d'une suite géométrique est :

\[ u_n = u_0 \times q^n \]

On a :

\begin{align*}
u_3 &= u_0 \times q^3 = 27 \\
u_6 &= u_0 \times q^6 = 729
\end{align*}

En divisant les deux équations :

\[ \dfrac{u_6}{u_3} = \dfrac{u_0 q^6}{u_0 q^3} = q^3 = \dfrac{729}{27} = 27 \]

Donc :

\[ q^3 = 27 \]

\[ q = \sqrt[3]{27} = 3 \]

Ensuite, on calcule \( u_0 \) :

\[ u_3 = u_0 \times q^3 \]

\[ u_0 = \dfrac{u_3}{q^3} = \dfrac{27}{3^3} = \dfrac{27}{27} = 1 \]

\end{enumerate}

\bigskip

\subsection*{Exercice 4 (6 points)}

\textbf{1. En numérotant les étages de haut en bas à partir de \( 0 \), soit \( u_n \) le nombre de cartes à l'étage \( n \), avec \( u_0 = 3 \).}

\begin{enumerate}[label=\alph*)]

\item \textbf{Donner \( u_1 \) et \( u_2 \).}

\textbf{Correction :}

Chaque étage comporte \( 3 \times (n + 1) \) cartes.

- Pour \( n = 1 \) : \( u_1 = 3 \times (1 + 1) = 3 \times 2 = 6 \)
- Pour \( n = 2 \) : \( u_2 = 3 \times (2 + 1) = 3 \times 3 = 9 \)

\bigskip

\item \textbf{Déterminer l'expression de \( u_n \).}

\textbf{Correction :}

\[ u_n = 3 \times (n + 1) \]

\end{enumerate}

\bigskip

\textbf{2. On cherche à calculer le nombre total de cartes \( s_n \) présentes dans un château ayant \( n \) étages.}

\begin{enumerate}[label=\alph*)]

\item \textbf{Donner \( s_0 \), \( s_1 \) et \( s_2 \).}

\textbf{Correction :}

\begin{itemize}
    \item \( s_0 = u_0 = 3 \)
    \item \( s_1 = u_0 + u_1 = 3 + 6 = 9 \)
    \item \( s_2 = u_0 + u_1 + u_2 = 3 + 6 + 9 = 18 \)
\end{itemize}

\bigskip

\item \textbf{Déterminer l'expression de \( s_n \).}

\textbf{Correction :}

\[ s_n = \sum_{k=0}^{n} u_k = \sum_{k=0}^{n} 3(k + 1) = 3 \sum_{k=0}^{n} (k + 1) \]

\[ \sum_{k=0}^{n} (k + 1) = \sum_{k=1}^{n+1} k = \dfrac{(n + 1)(n + 2)}{2} \]

Donc :

\[ s_n = 3 \times \dfrac{(n + 1)(n + 2)}{2} = \dfrac{3(n + 1)(n + 2)}{2} \]

\end{enumerate}

\bigskip

\textbf{3. On dispose de 315 cartes. Combien d'étages peut-on construire ?}

\textbf{Correction :}

On cherche le plus grand \( n \) tel que :

\[ s_n \leq 315 \]

\[ \dfrac{3(n + 1)(n + 2)}{2} \leq 315 \]

\[ (n + 1)(n + 2) \leq \dfrac{315 \times 2}{3} = 210 \]

Donc :

\[ (n + 1)(n + 2) \leq 210 \]

Testons des valeurs entières pour \( n \) :

\begin{itemize}
    \item Pour \( n = 12 \) :

    \[ (12 + 1)(12 + 2) = 13 \times 14 = 182 \leq 210 \]

    \item Pour \( n = 13 \) :

    \[ (13 + 1)(13 + 2) = 14 \times 15 = 210 \]

    \item Pour \( n = 14 \) :

    \[ (14 + 1)(14 + 2) = 15 \times 16 = 240 > 210 \]
\end{itemize}

Donc, le nombre maximal d'étages est \( n = 13 \).

\textbf{Réponse :} On peut construire \textbf{13 étages} avec 315 cartes.

\end{document}
